\section{Nested Loops: the Inner Loop as a Value}\label{sec:nested}

Four nested loops close with the database alone: an inner counter (\texttt{nested\_counter}), its final values
summed (\texttt{nested\_sum}), a variable incremented by the inner loop (\texttt{nested\_mult}), and the sums of
the rows of a matrix stored by the outer loop (\texttt{row\_sums}). A fifth, the transpose, needs one lemma
that states what the two loops together compute. This section shows what the closed form of nested loops of
Section~\ref{sec:teq} provides and what it leaves to a lemma.

\subsection{The nested sum}
\begin{lstlisting}
method main(m: int, n: int) returns (c: int)
  requires 0 <= m && 0 <= n
  ensures c == m * n
{
  var i: int; var j: int;
  c := 0; i := 0;
  while i < m {
    j := 0;
    while j < n { j := j + 1; }
    c := c + j;
    i := i + 1;
  }
}
\end{lstlisting}
With $J_7 = \cur(\tr{0}{\bl{\cnt(0, 1, n_7)}}{j, k \mapsto j + 1})$ for the final value of $j$ in an outer
iteration with $n_7$ inner iterations, the system of the outer loop is
\begin{align}
  K &= \cnt(0, 1, n') \qquad c = \cur(C) \qquad i = \cur(I) \nonumber\\
  C &= \tr{0}{\bl{K}, \seq{N}}{c, k, n_7 \mapsto \upd{\neg(J_7 < n)}{c + J_7}} \label{eq:ns1}\\
  I &= \tr{0}{\bl{K}, \seq{N}}{i, k, n_7 \mapsto \upd{\neg(J_7 < n)}{i + 1}} \nonumber\\
  &\mathsf{invs}(\mathsf{for}\ \bl{K}, \bl{C}, \bl{I}, \seq{N}\ \mathsf{check}\ k, c, i, n_7 \mapsto i < m \nonumber\\
  &\qquad \wedge \al{\bl{\cnt(0, 1, n_7)}, \bl{\tr{0}{\ldots}{\ldots}}}{k', j \mapsto j < n} \wedge \neg(J_7 < n)) \label{eq:ns2}
\end{align}
and the exit $\neg(i < m)$. The update has one case, guarded by the exit of the inner loop: the path through the
body passes the inner loop only once it has stopped. The inner loop occurs as a term in the updates and in the
constraint; only the numbers $\seq{N}$ of its iterations are fresh.

\paragraph{What the proof needs.} Everything is a lemma on single loops, applied inside the outer combinators:
\begin{enumerate}[nosep]
\item The inner trace is a counter (\lemma{counter\_def}), so $J_7 = \intf(n_7)$ (\lemma{current\_counter}).
  This happens inside the update and inside the constraint alike, since the inner loop is a term like any
  other.
\item Inside the constraint \eqref{eq:ns2}, the inner test over the counter becomes $\intf(n_7) \le n$
  (\lemma{invs\_counter\_upper\_nonnegative}), and with the inner exit $\neg(\intf(n_7) < n)$ the number of
  inner iterations is determined: $\intf(n_7) = n$ in every outer iteration.
\item A choice whose values the constraint determines is a constant list (\lemma{consts\_choice}):
  $\seq{N} = \bl{\mathsf{consts}(n', \natf(n))}$. Its reads in the updates become $\natf(n)$
  (\lemma{earlier\_consts}), and the guard $\neg(J_7 < n)$ becomes $\mathsf{true}$.
\item The trace of $c$ adds $n$ in every iteration: it is folded into $\sums$ over a constant list
  (\lemma{sums}), whose current value is $n \cdot \intf(n')$ (\lemma{current\_sums}, \lemma{sum\_consts}). The
  outer constraint and exit give $\intf(n') = m$, and the arithmetic substitutes the solution into the
  polynomial: $c = m \cdot n$.
\end{enumerate}
\texttt{nested\_counter} stops after the second step, and \texttt{nested\_mult} differs in the third: the inner
loop increments the outer variable itself, so its trace is a counter from the previous value $r$, whose final
value is $r + \intf(n_7)$, and the outer trace of $r$ adds $n$ per iteration as above.

\paragraph{Why the inner loop must be a value.} If the inner loop were described by fresh lists of the outer
loop, one per inner trace and its final value, each with its own iteration constraint, then the facts about
one of them (the inner counter is $\cnt(0, 1, n_7)$) would sit in a different combinator than the facts that
need them (the final value is the current value of that counter), and every substitution between them would
be a lemma of its own, one per kind of inner loop. As a value of the update, the inner loop is one term, and
the lemmas on single loops apply to it where it occurs, with the facts of the outer iteration in the context.
The only thing that cannot be a value is the number of inner iterations, which is what the inner exit
determines: it stays a choice, and \lemma{consts\_choice} resolves it as soon as the constraint fixes it.

\subsection{The row sums}
\begin{lstlisting}
method main(A: array<int, array<int, int>>, m: int, n: int) returns (s: array<int, int>)
  requires 0 <= m && 0 <= n
  ensures forall i: int :: 0 <= i < m ==>
    s[i] == sum(selects(A[i], ⌊counter(0, 1, natural(n))⌋))
{
  var i: int; var j: int; var t: int;
  i := 0;
  while i < m {
    t := 0; j := 0;
    while j < n { t := t + A[i][j]; j := j + 1; }
    s := s[i := t];
    i := i + 1;
  }
}
\end{lstlisting}
The inner loop sums row $i$, the outer loop stores the sum at position $i$. The update of the trace of $s$ is
\begin{align*}
  s[i := \cur(\mathsf{trace}(0,\ \mathsf{for}\ &\bl{\cnt(0, 1, n_{10})}, \bl{\tr{0}{\bl{\cnt(0, 1, n_{10})}}{j, k \mapsto j + 1}}\\
  &\mathsf{do}\ t, k, j \mapsto t + A[i][j]))],
\end{align*}
with the inner trace of $j$ among the dependencies of the inner trace of $t$. The numbers of inner iterations
are resolved to $\natf(n)$ as in the nested sum. Then the database does the rest: the inner trace of $j$ is a
counter, the reads $A[i][j]$ over it are $\mathsf{selects}(A[i], \bl{\cnt(0, 1, \natf(n))})$ (\lemma{selects}),
the inner trace of $t$ is a $\sums$ over them, and its current value is their sum (\lemma{current\_sums}). The
outer trace stores these sums at the positions of a counter (\lemma{stores}), and the read $s[i]$ of the
postcondition is the stored value at $i$, by \lemma{current\_stores} with the membership and the position of
$i$ in the counter (\lemma{occurs\_counter\_up}, \lemma{indexof\_counter\_up}): exactly the sum that the
postcondition states. Nothing in this proof is about nested loops; it is the lemmas of single loops, applied
in the order in which the terms appear.

\subsection{The transpose: one lemma for the nest}\label{sec:transpose}
\begin{lstlisting}
lemma {:rw} current_transpose(B0: array<int, array<int, int>>, A: array<int, array<int, int>>,
                              m: int, n: int, i: int, j: int)
  requires 0 <= i < m && 0 <= j < n
  ensures current(trace(B0, for ⌊counter(0, 1, natural(m))⌋ do B, k ::
      current(trace(B, for ⌊counter(0, 1, natural(n))⌋ do Bi, ki ::
        Bi[ki := Bi[ki][k := A[k][ki]]]))))[j][i]
    == A[i][j]

method main(A: array<int, array<int, int>>, m: int, n: int)
  returns (B: array<int, array<int, int>>)
  requires 0 <= m && 0 <= n
  ensures forall i: int, j: int :: 0 <= i < m && 0 <= j < n ==> B[j][i] == A[i][j]
{
  var i: int; var j: int;
  i := 0;
  while i < m {
    j := 0;
    while j < n { B := B[j := B[j][i := A[i][j]]]; j := j + 1; }
    i := i + 1;
  }
}
\end{lstlisting}
The inner loop stores column $i$ of $A$ into row positions of $B$, a store into an array that is itself an
element of $B$. No lemma of the database describes two-dimensional stores, so the program states one: the left-hand
side is the nest as the system has it once the numbers of inner iterations are resolved to $\natf(n)$, the outer
trace whose update is the current value of the inner trace, and the right-hand side the element of the transpose.
The lemma is a summary of the whole nest in one line, with the side conditions that make the read meaningful.
The program is otherwise proved by the database, as the row sums: the counters, the choice, and the read of the
postcondition at $j$, $i$ within the bounds.

\paragraph{What it needs.} The lemma matches because its left-hand side is the term of the goal: an outer trace
over a counter whose update contains the inner trace as a value. A summary of the inner loop alone would have
been a rule on $\cur(\tr{B}{\bl{\cnt(0, 1, \natf(n))}}{\ldots})$ with the column as a named function, followed by
a rule on the outer trace of these columns; the closed form lets one lemma state the result of both, at the
price of a lemma specific to the nest. The same file in the folder of the algorithms keeps two folds, of the
columns and of the outer results, and four lemmas written for inner loops described by fresh lists; they do not
match the closed form, and one goal remains (Section~\ref{sec:open}).
